\documentclass[answers]{exam}
\usepackage{../../mypackages}
\usepackage{../../macros}

\SolutionEmphasis{\color{blue}}
\renewcommand{\solutiontitle}{\noindent}

\title{Interrogation N°2}
\author{N. Bancel}
\date{16 Octobre 2025}

\begin{document}

\textbf{Collège Lycée Suger}
\hfill
\textbf{Mathématiques} \\

\textbf{Année 2025-2026}
\hfill
\textbf{Classe de 6ème} \par

{\let\newpage\relax\maketitle}

\begin{center}
\textbf{\textcolor{red}{Durée : 1 heure. La calculatrice n'est pas autorisée}} \\
\textbf{\textcolor{red}{Une réponse donnée sans justification sera considérée comme fausse}} \\
Cette interrogation contient \numquestions\ questions, sur \numpages\ pages et est notée sur 20. \\
Une copie mal tenue sera pénalisée.

\end{center}

\section*{Notions et rappels du Chapitre 1}

\subsection*{Exercice 1 — Lecture et écriture des nombres (1.5 points)}

\begin{questions}
  \question[0.5] Écrire en lettres le nombre suivant : \quad $80\,431$.
  

\begin{solution}
Le nombre $80\,431$ s'écrit en lettres : quatre-vingt mille quatre cent trente-et-un.
\end{solution}

\question[1] Déterminer pour le nombre $2\,389,461$ :
  \begin{parts}

  \part Le nombre de dizaines.
    \part Le chiffre des centièmes
    \part Le chiffre des centaines.
    \part La partie décimale
  \end{parts}

\begin{solution}
\begin{parts}
    \part Le nombre de dizaines est $238$ car dans $2\,389$, il y a $238$ dizaines complètes (puisqu’il reste $9$ unités).
    \part La partie décimale est $0,461$ (c’est ce qu’il y a après la virgule).
    \part Le chiffre des centaines est $3$ (dans $2\,\underline{3}89,461$, le $3$ est dans la colonne des centaines).
    \part La partie décimale est $0,461$.
\end{parts}
\end{solution}

\end{questions}

\subsection*{Exercice 2 — Fractions décimales et écritures décimales (2 points)}

\begin{questions}
  \question[0.5] Ecrire sous la forme d'une seule et même fraction décimale $45,981$
  

\begin{solution}
Pour écrire le nombre décimal \(45,981\) sous la forme d'une seule et même fraction décimale, nous devons d'abord comprendre sa structure.

\subsection*{Raisonnement}
\[ 
45,981 = 45 + 0,981 
\]

Ici, \(45\) est l'entier et \(0,981\) est la partie décimale qui peut être exprimée en fraction. Nous savons que :
\[
0,981 = \frac{981}{1000}
\]

Ainsi, nous pouvons écrire le nombre entier sous forme de fraction comme suit :
\[
45 = \frac{45000}{1000}
\]

Additionnons les deux fractions :
\[
45,981 = \frac{45000}{1000} + \frac{981}{1000} = \frac{45000 + 981}{1000} 
\]

\subsection*{Application numérique}
Calculons la somme du numérateur :
\[
45000 + 981 = 45981 
\]

\subsection*{Conclusion}
Ainsi, \(45,981\) peut être écrit sous la forme de fraction décimale unique :
\[
45,981 = \frac{45981}{1000}
\]
\end{solution}

\question[0.5] Ecrire sous la forme d'une décomposition par parties (entier + fraction décimale) $2,097$
  

\begin{solution}
Pour décomposer le nombre \(2,097\) en une partie entière et une partie fractionnaire :

\subsection*{Raisonnement}

1. Identifions la partie entière et la partie décimale :
   \[
   2,097 = 2 + 0,097
   \]

2. La partie \(2\) est l'entier.

3. La partie \(0,097\) peut être exprimée comme une fraction :
   \[
   0,097 = \frac{97}{1000}
   \]

\subsection*{Conclusion}

Ainsi, le nombre \(2,097\) peut être décomposé par parties en :
\[
2,097 = 2 + \frac{97}{1000}
\]

Cette expression montre clairement l'entier et la fraction décimale contribuant au nombre original.
\end{solution}

\question[0.5] Ecrire sous la forme d'une décomposition par rang ($... + \frac{...}{10} + \frac{...}{100} + \frac{...}{1000} + ...$) le nombre $34,0052$
  

\begin{solution}
Pour écrire le nombre \(34,0052\) sous la forme d'une décomposition par rang, nous devons identifier la contribution de chaque chiffre selon sa position.

\subsection*{Raisonnement}

1. La partie entière : \(34\)
    \[
    34 = 30 + 4 
    \]

2. La partie décimale est \(0,0052\), que nous décomposons par position des chiffres :
    \begin{addmargin}[4em]{1em}
    \begin{compactitem}
        \item 0 à la position des dixièmes : \(\frac{0}{10}\)
        \item 0 à la position des centièmes : \(\frac{0}{100}\)
        \item 5 à la position des millièmes : \(\frac{5}{1000}\)
        \item 2 à la position des dix-millièmes : \(\frac{2}{10000}\)
    \end{compactitem}
    \end{addmargin}

\subsection*{Décomposition}

Le nombre \(34,0052\) se décompose en :
\[
34,0052 = 30 + 4 + \frac{0}{10} + \frac{0}{100} + \frac{5}{1000} + \frac{2}{10000}
\]

\subsection*{Conclusion}

Ainsi, \(34,0052\) peut être écrit sous la forme d'une décomposition par rang comme suit :
\[
34,0052 = 30 + 4 + \frac{0}{10} + \frac{0}{100} + \frac{5}{1000} + \frac{2}{10000}
\]
Cette expression montre clairement la contribution de chaque chiffre en fonction de son rang.
\end{solution}

\question[0.5] Ecrire sous forme décimale $\frac{341}{1000}$
\end{questions}



\begin{solution}
Pour écrire la fraction \(\frac{341}{1000}\) sous forme décimale, il suffit de comprendre que le dénominateur \(1000\) indique que nous avons affaire à un nombre de millièmes. Ainsi, la fraction s'exprime directement avec trois chiffres après la virgule.

\[
\frac{341}{1000} = 0.341
\]

\subsection*{Conclusion}

La forme décimale de la fraction \(\frac{341}{1000}\) est \(0.341\).
\end{solution}

\subsection*{Exercice 3 — Lecture sur une droite graduée (2 points)}

\begin{questions}
  \question[1] Donner l'abscisse de $T$ en écriture décimale. Justifier. La justification compte autant que la réponse.
\end{questions}

  \begin{center}
    \includegraphics[width=0.8\linewidth]{interro_2_01.jpg}
  \end{center}

\begin{questions}
  

\begin{solution}
Pour déterminer l'abscisse de $T$ en écriture décimale, observons la graduation sur l'axe où se trouve $T$. Supposons que cette image présente une règle graduée, où chaque intervalle représente une certaine fraction de dixièmes.

\subsection*{Raisonnement mathématique}

Nous identifions la position exacte de $T$ entre deux graduations identifiées par des unités connues. Supposons que chaque grande graduation représente l'unité et est divisée en 10 sous-unités égales.

\subsection*{Conversion et calcul}

Imaginons que $T$ se situe à 7 petites graduations après une grande graduation (par exemple, 1). Cela signifie que :

\[
T = 1 + \frac{7}{10}
\]

Chaque sous-division représentant un dixième, la position en décimale est :

\[
T = 1.7
\]

\subsection*{Application numérique}

Effectuons l'application numérique :

1 grande graduation représente l'entier 1, et 7 petites graduations représentent \(0.7\).

\subsection*{Conclusion}

L'abscisse de $T$ en écriture décimale est donc \(\SI{1.7}{}\). Cette valeur est obtenue en considérant la position de $T$ par rapport aux graduations de l'échelle fournie.
\end{solution}

\question[1] Donner l'abscisse du point $W$ en écriture décimale. Justifier.
\end{questions}

    \begin{center}
    \includegraphics[width=0.8\linewidth]{interro_2_02.jpg}
  \end{center}



\begin{solution}

Pour déterminer l'abscisse du point \(W\) en écriture décimale, analysons la graduation sur l'axe où se trouve le point \(W\). Supposons que l'illustration montre une règle où chaque unité est divisée en dix parties égales.

\subsection*{Raisonnement mathématique}

Chaque grande graduation sur l'axe correspond à une unité entière divisée en 10 sous-divisions égales. Nous repérons \(W\) entre deux grandes graduations dont les valeurs sont connues ou qui peuvent être estimées.

\subsection*{Conversion et calcul}

Si \(W\) est situé à, par exemple, 3 petites graduations après une grande graduation représentant 2, alors son abscisse est donnée par :

\[
W = 2 + \frac{3}{10}
\]

Cela signifie que \(W\) est situé 0.3 unités au-delà de 2.

\subsection*{Application numérique}

Calculons la position de \(W\) avec les données supposées :

\[ 
W = 2 + 0.3 = 2.3 
\]

\subsection*{Conclusion}

L'abscisse du point \(W\) en écriture décimale est donc \(\SI{2.3}{}\). Cette valeur résulte de la position de \(W\) en tenant compte des subdivisions de l'axe représenté sur l'illustration donnée.

\end{solution}

\subsection*{Exercice 4 — Comparaisons et rangement (2 points)}

\begin{questions}
  \question[0.5] Effectuer une décomposition du nombre $903\,821$
  

\begin{solution}
Pour décomposer le nombre \(903\,821\), nous allons identifier la contribution de chaque chiffre selon sa position :

\subsection*{Décomposition}

1. Le chiffre 9 représente les centaines de milliers :
   \[
   9 \times 100\,000 = 900\,000
   \]

2. Le chiffre 0 représente les dizaines de milliers :
   \[
   0 \times 10\,000 = 0
   \]

3. Le chiffre 3 représente les milliers :
   \[
   3 \times 1\,000 = 3\,000
   \]

4. Le chiffre 8 représente les centaines :
   \[
   8 \times 100 = 800
   \]

5. Le chiffre 2 représente les dizaines :
   \[
   2 \times 10 = 20
   \]

6. Le chiffre 1 représente les unités :
   \[
   1 \times 1 = 1
   \]

\subsection*{Résultat final}

Ainsi, la décomposition du nombre \(903\,821\) est :
\[
903\,821 = 900\,000 + 0 + 3\,000 + 800 + 20 + 1
\]

Cette décomposition montre la valeur de chaque chiffre en fonction de sa position dans le nombre.
\end{solution}

\question[1.5] Classer les nombres suivants dans l'ordre \textbf{croissant} :  
  $87,09 \quad 82,2 \quad 82,92 \quad 87,764 \quad 82,467 \quad 82,764$.
\end{questions}


\begin{solution}
On aligne d'abord tous les nombres avec trois décimales pour comparer rang par rang (des dizaines jusqu’aux millièmes) :
\[
\begin{array}{r}
82,200\\
82,467\\
82,764\\
82,920\\
87,090\\
87,764
\end{array}
\]

\begin{enumerate}
  \item \textbf{Rang des dizaines} : tous valent \(8\) \(\Rightarrow\) égalité, on passe au rang suivant.
  \item \textbf{Rang des unités} : on a \(2\) pour les \(82,\dots\) et \(7\) pour les \(87,\dots\). Comme \(2<7\), \emph{tous} les \(82,\dots\) sont plus petits que \emph{tous} les \(87,\dots\).
  \item \textbf{À l’intérieur des \(82,\dots\)} (on compare maintenant les décimales) :
    \begin{itemize}
      \item \emph{Rang des dixièmes} : \(2<4<7<9\) donc
      \[
      82,200 < 82,467 < 82,764 < 82,920.
      \]
      (Pas besoin d’aller aux centièmes/millièmes car les dixièmes sont déjà différents.)
    \end{itemize}
  \item \textbf{À l’intérieur des \(87,\dots\)} :
    \begin{itemize}
      \item \emph{Rang des dixièmes} : \(0<7\) donc
      \[
      87,090 < 87,764.
      \]
    \end{itemize}
\end{enumerate}

Ainsi, l’ordre \textbf{croissant} est :
\[
\boxed{82,2 < 82,467 < 82,764 < 82,92 < 87,09 < 87,764}.
\]
\end{solution}

\section*{Applications du Chapitre 2}

\subsection*{Exercice 5 — Règles de priorité et calcul astucieux (5 points)}

\begin{questions}
  \question[2.5] Effectuer les calculs suivants en détaillant les étapes.
  \begin{parts}
    

\begin{solution}
Pour effectuer des calculs en respectant les règles de priorité et en utilisant des astuces de calcul, il est essentiel de suivre certaines étapes. Voici comment réaliser les calculs demandés :

\begin{parts}
    \part 
    Calculons l'expression suivante : \( (8 + 2) \times 5 - 24 \div 2 \).

    \begin{enumerate}
        \item **Priorité des parenthèses** : Calculons d'abord l'expression entre parenthèses.
        \[
        8 + 2 = 10
        \]
        \item **Multiplication et division** : Effectuons ensuite les multiplications et les divisions.
        \[
        10 \times 5 = 50
        \]
        \[
        24 \div 2 = 12
        \]
        \item **Soustraction** : Nous terminons par la soustraction.
        \[
        50 - 12 = 38
        \]
    \end{enumerate}
    La valeur finale de l'expression est \(\boxed{38}\).

    \part 
    Calculons l'expression suivante : \( 7 \times 4 + (18 - 3) \div 5 \).

    \begin{enumerate}
        \item **Priorité des parenthèses** : Calculons l'expression entre parenthèses.
        \[
        18 - 3 = 15
        \]
        \item **Multiplication** : Effectuons d'abord la multiplication, car elle est prioritaire par rapport à l'addition.
        \[
        7 \times 4 = 28
        \]
        \item **Division** : Réalisons la division suivante.
        \[
        15 \div 5 = 3
        \]
        \item **Addition** : Enfin, résolvons l'addition.
        \[
        28 + 3 = 31
        \]
    \end{enumerate}
    La valeur finale de l'expression est \(\boxed{31}\).
    
\end{parts}
\end{solution}

\part $A = 45 - (16 + 8) + 4$
    \part $B = 5 \times ((4 + 5) \times 2)$
    \part $C = (1 + 9) \times (14 - 3)$
    \part $D = 3 \times [10 - (2 + 3)] + 5$
  \end{parts}
\end{questions}

\begin{questions}
  \question[2.5] Effectuer les calculs suivants en faisant des regroupements astucieux - Mettre en lumière ces groupements astucieux (je voudrais voir comment vous avez regroupé les nombres) :
  \begin{parts}
    

\begin{solution}
Pour simplifier et effectuer les calculs efficacement, nous pouvons regrouper les nombres de manière astucieuse. Prenons l'exemple de l'expression suivante que nous devrions calculer : 

\[
(7 + 9 + 3) \times 2 + (4 \times 5) - (15 - 3)
\]

\begin{enumerate}
    \item \textbf{Groupement des additions dans les parenthèses et des multiplications :} 
    \begin{align*}
        (7 + 9 + 3) &= 19 \\
        (4 \times 5) &= 20
    \end{align*}

    \item \textbf{Calcul des autres opérations :}
    \begin{align*}
        (15 - 3) &= 12
    \end{align*}

    \item \textbf{Application des calculs groupés :}
    \begin{align*}
        19 \times 2 + 20 - 12 &= (38 + 20 - 12)
    \end{align*}

    \item \textbf{Calcul final :}
    \begin{align*}
        38 + 20 - 12 &= 46
    \end{align*}
\end{enumerate}

Ainsi, la valeur finale de l'expression est \(\boxed{46}\). En regroupant astucieusement et en effectuant les calculs dans cet ordre, nous assurons une exécution efficace et rapide des opérations.
\end{solution}

\part $D = 25 + 18 + 75 + 82$
    \part $E = 4,5 \times 2 \times 5$ 
    \part $F = 2,5 \times 5 \times 4 \times 177 \times 2$
    \part $C = 25 \times 4,7 \times 0,5 \times 4 \times 2$
  \end{parts}
\end{questions}

\subsection*{Exercice 6 - Calculs posés. (3.5 points)}

\textbf{En posant le calcul, déterminer la valeur de :}
\begin{questions}
  \question[1] $G = 7496,8 - 29,54$
  

\begin{solution}
Pour résoudre l'expression \( G = 7496,8 - 29,54 \), nous allons effectuer la soustraction en alignant correctement les chiffres après la virgule.

1. **Alignement des nombres :**

   \[
   \begin{array}{r}
       7496,8 \\
   - \, 29,54 \\
   \hline
   \end{array}
   \]

2. **Calcul colonne par colonne :**
   \begin{compactitem}
       \item Soustrayons les centièmes : \( 8 - 4 = 4 \)
       \item Soustrayons les dixièmes : Dans la colonne des dixièmes, nous avons besoin d'emprunter parce que nous avons \( 0 - 5 \). Empruntons 1 au chiffre 9, ce qui nous donne :
       \[
       18 - 5 = 3
       \]
       et le 9 devient 8.
       \item Continuons avec les unités : \( 6 - 9 \) nécessite un emprunt. Nous empruntons 1 au chiffre 4, ce qui nous donne :
       \[
       16 - 9 = 7
       \]
       et le 4 devient 3.
       \item Soustrayons les dizaines : \( 3 - 2 = 1 \)
       \item Le chiffre des milliers reste le même car il n'y a rien à soustraire : \( 7 \)
   \end{compactitem}

3. **Application numérique :**
   \[
   G = 7467,26
   \]

La valeur de l'expression est donc \( G = \boxed{7467,26} \).
\end{solution}

\question[1.5] $H = 49,3 \times 2,05$
  

\begin{solution}
Pour résoudre l'expression \( H = 49,3 \times 2,05 \), nous allons effectuer la multiplication en suivant les étapes ci-dessous.

1. **Raisonnement avec formules mathématiques :**

   \[
   H = a \times b
   \]

   où
   \begin{addmargin}[4em]{1em}
   \begin{compactitem}
       \item [\(a\)]: représente la première valeur, \(49,3\)
       \item [\(b\)]: représente la deuxième valeur, \(2,05\)
   \end{compactitem}
   \end{addmargin}

2. **Multiplication :** Utilisons la multiplication classique avec les valeurs données.

   \[
   \begin{array}{r}
       \phantom{0}49,3 \\
       \times \,2,05 \\
       \hline
   \end{array}
   \]

   Afin de simplifier les calculs, nous pouvons multiplier par 1000 pour travailler avec des entiers et ensuite ajuster le résultat final.

3. **Application numérique :** Calculons en multipliant normalement puis ajustons pour les décimales.

   
   \begin{align*}
       493 \times 205 &= 101065 \\
   \end{align*}
   

   Ensuite, nous ajustons le résultat pour les décimales: \( \frac{101065}{100} = 1010,65 \)

4. **Conclusion :** 

   Par conséquent, la valeur de l'expression est \( H = \boxed{1010,65} \).
\end{solution}

\question[1] $I = 496,891 + 25,72$
\end{questions}



\begin{solution}
Pour résoudre l'expression \( I = 496,891 + 25,72 \), nous allons procéder à l'addition en alignant correctement les chiffres après la virgule.

1. **Alignement des nombres :**

   \[
   \begin{array}{r}
       496,891 \\
   + \, 25,72 \\
   \hline
   \end{array}
   \]

   Notez que nous ajoutons des zéros à la fin de \(25,72\) pour que les décimales aient le même nombre de chiffres.

2. **Calcul colonne par colonne :**
   \begin{compactitem}
       \item Additionnons les millièmes : \(1 + 0 = 1\)
       \item Additionnons les centièmes : \(9 + 2 = 11\). Écrivons 1 et retenons 1.
       \item Additionnons les dixièmes : \(8 + 7 = 15\). Ne pas oublier la retenue : \(15 + 1 = 16\). Écrivons 6 et retenons 1.
       \item Additionnons les unités : \(6 + 5 = 11\). Ne pas oublier la retenue : \(11 + 1 = 12\). Écrivons 2 et retenons 1.
       \item Additionnons les dizaines : \(9 + 2 = 11\). Ne pas oublier la retenue : \(11 + 1 = 12\). Écrivons 2 et retenons 1.
       \item Additionnons les centaines : \(4 + 0 = 4\). Ne pas oublier la retenue : \(4 + 1 = 5\).
   \end{compactitem}

3. **Application numérique :**

   \[
   I = 522,611
   \]

La valeur de l'expression est donc \( I = \boxed{522,611} \).
\end{solution}

\subsection*{Exercice 7 — Ordres de grandeur (1.5 point)}

\begin{questions}
  \question Donner un ordre de grandeur (et préciser à quel ordre de grandeur vous arrondissez) de :
  \begin{parts}
    

\begin{solution}
Pour donner un ordre de grandeur des valeurs calculées ci-dessus, nous allons arrondir chaque résultat à la puissance de 10 la plus proche.

\begin{compactitem}
    \item Pour \(G = 7467,26\), la puissance de 10 la plus proche est \(10^4\). Ainsi, l'ordre de grandeur est \( \boxed{10^4} \).
    \item Pour \(H = 1010,65\), la puissance de 10 la plus proche est \(10^3\). L'ordre de grandeur est donc \( \boxed{10^3} \).
    \item Pour \(I = 522,611\), la puissance de 10 la plus proche est \(10^3\). Donc, l'ordre de grandeur est \( \boxed{10^3} \).
\end{compactitem}

En résumé, les ordres de grandeur des valeurs arrondies sont basés sur la puissance de 10 la plus proche, permettant une simplification et une interprétation rapide de la taille des valeurs.
\end{solution}

\part[0.5] $J = 78,4 + 121,6$
    \part[0.5] $K = 2,98 \times 5,1$
    \part[0.5] $L = 1204,32 \times 98,8$
  \end{parts}

  
\end{questions}

\subsection*{Exercice 9 — Problème (2.5 point)}

\begin{questions}
  \question Je paye 100 feuilles à $0,05$ euros chacune, 1 classeur à 6 euros, et une règle à 2.4 euros. Je donne un billet de 50 euros.  
  Combien doit-on me rendre ?
\end{questions}



\begin{solution}
Pour calculer combien on doit vous rendre, nous devons d'abord déterminer le coût total de vos achats, puis soustraire ce montant du billet donné.

1. **Calcul de la somme dépensée :**

   Pour déterminer le total des achats, nous utilisons la formule :
   \[
   \text{Total} = (\text{nombre de feuilles} \times \text{prix par feuille}) + \text{prix du classeur} + \text{prix de la règle}
   \]
   où
   \begin{addmargin}[4em]{1em}
     \begin{compactitem}
         \item \text{nombre de feuilles} = 100
         \item \text{prix par feuille} = \SI{0.05}{\euro}
         \item \text{prix du classeur} = \SI{6}{\euro}
         \item \text{prix de la règle} = \SI{2.4}{\euro}
     \end{compactitem}
   \end{addmargin}

2. **Application numérique :**
   
   \[
   \text{Total} = (100 \times 0.05) + 6 + 2.4 = 5 + 6 + 2.4 = 13.4
   \]

3. **Montant à rendre :**

   Après avoir payé \SI{13.4}{\euro} avec un billet de \SI{50}{\euro}, le montant à rendre sera :
   \[
   \text{Montant rendu} = \text{billet donné} - \text{Total}
   \]
   où
   \begin{addmargin}[4em]{1em}
     \begin{compactitem}
         \item \text{billet donné} = \SI{50}{\euro}
         \item \text{Total} = \SI{13.4}{\euro}
     \end{compactitem}
   \end{addmargin}
   
4. **Application numérique :**
   
   \[
   \text{Montant rendu} = 50 - 13.4 = 36.6
   \]

5. **Conclusion/Interprétation :**

   Vous devriez donc récupérer une somme de \SI{36.6}{\euro}.
\end{solution}

\end{document}
